\documentclass[11pt]{article}
\usepackage[margin=1in]{geometry}
\usepackage{amsmath,amssymb}
\usepackage{booktabs}
\usepackage{hyperref}
\usepackage{enumitem}

\title{Methodology Summary: Epidemiological Parameter Estimation\\for U.S. Counties (March 2020 -- December 2021)}
\author{SIHRS Model Parameter Analysis}
\date{}

\begin{document}
\maketitle

\section{Overview}

This document describes the methodology used to calculate three key epidemiological parameters for all U.S. counties during the COVID-19 pandemic (March 2020 -- December 31, 2021):

\begin{enumerate}
    \item \textbf{$p_{IH}$}: Probability of hospitalization given infection
    \item \textbf{$p_{ID}$}: Probability of death given infection
    \item \textbf{CFR}: Case Fatality Rate (cumulative deaths / cumulative cases)
\end{enumerate}

\section{Data Sources}

\subsection{Hospitalization Data}
\begin{itemize}
    \item \textbf{Source}: HealthData.gov COVID-19 Reported Patient Impact and Hospital Capacity by Facility
    \item \textbf{URL}: \url{https://healthdata.gov/Hospital/COVID-19-Reported-Patient-Impact-and-Hospital-Capa/anag-cw7u}
    \item \textbf{File}: \texttt{covid\_hospital\_report.csv} (679 MB)
    \item \textbf{Temporal Resolution}: Weekly (7-day averages)
    \item \textbf{Key Column}: \texttt{total\_adult\_patients\_hospitalized\_confirmed\_and\_suspected\_covid\_7\_day\_avg}
\end{itemize}

\subsection{Case and Death Data}
\begin{itemize}
    \item \textbf{Source}: The New York Times COVID-19 Data Repository
    \item \textbf{URL}: \url{https://github.com/nytimes/covid-19-data}
    \item \textbf{File}: \texttt{nyt\_covid\_data.csv} (100 MB)
    \item \textbf{Temporal Resolution}: Daily
    \item \textbf{Key Columns}: \texttt{date}, \texttt{county}, \texttt{state}, \texttt{fips}, \texttt{cases}, \texttt{deaths}
\end{itemize}

\section{Data Extraction and Processing}

\subsection{Step 1: Hospitalization Data Extraction}

The raw hospitalization data from HealthData.gov contains facility-level reports. We processed this data as follows:

\begin{enumerate}
    \item \textbf{Filtering}: Selected records within the date range (August 2, 2020 -- December 26, 2021)
    \item \textbf{Quality Control}: Excluded records with null or invalid values ($-999999$)
    \item \textbf{County Identification}: Used 5-digit FIPS codes for county matching
    \item \textbf{Aggregation}: Summed hospitalization values across multiple facilities within the same county on the same date
    \item \textbf{Output}: Created individual CSV files for 2,446 counties, organized by state:
    \begin{verbatim}
    CaseStudies/hospitalization_data/by_state/<STATE>/<county>_<fips>.csv
    \end{verbatim}
\end{enumerate}

\subsection{Step 2: Active Cases Calculation}

Active cases represent the number of currently infected individuals at any given time. We calculated active cases using a 14-day recovery window:

\begin{equation}
    \text{Active Cases}(t) = \max\left(\text{Cases}(t) - \text{Cases}(t-14), 0\right)
\end{equation}

This approach assumes:
\begin{itemize}
    \item Average recovery time of 14 days
    \item Recovered individuals are no longer counted as active cases
    \item Negative values are clipped to zero
\end{itemize}

\subsection{Step 3: Daily Deaths Calculation}

Daily deaths were derived from cumulative death counts:

\begin{equation}
    \text{Daily Deaths}(t) = \text{Deaths}(t) - \text{Deaths}(t-1)
\end{equation}

\section{Parameter Calculations}

\subsection{$p_{IH}$: Probability of Hospitalization Given Infection}

\subsubsection{Formula}
\begin{equation}
    p_{IH} = \frac{\text{Hospitalized Patients}(T)}{\text{Active Cases}(T-14)}
\end{equation}

\subsubsection{Methodology}
\begin{enumerate}
    \item \textbf{14-day lag}: Accounts for the time from infection to hospitalization:
    \begin{itemize}
        \item Days 0--5: Incubation period (asymptomatic)
        \item Days 5--10: Symptom development
        \item Days 10--14: Hospitalization decision for severe cases
    \end{itemize}
    \item \textbf{Weekly data handling}: Used 2-row lookback in weekly hospitalization data (approximately 14 days)
    \item \textbf{Date validation}: Only included calculations where the actual date difference was 10--15 days
    \item \textbf{Matching tolerance}: $\pm 1$ day tolerance for matching active cases dates
    \item \textbf{Inclusion of zeros}: Calculations where hospitalized = 0 were included to avoid upward bias
\end{enumerate}

\subsubsection{Output}
\begin{itemize}
    \item Counties analyzed: 2,438
    \item Total calculations: 170,487
    \item Output files: \texttt{CaseStudies/p\_IH\_analysis/}
\end{itemize}

\subsection{$p_{ID}$: Probability of Death Given Infection}

\subsubsection{Formula}
\begin{equation}
    p_{ID} = \frac{\text{Daily Deaths}(T)}{\text{Active Cases}(T-20)}
\end{equation}

\subsubsection{Methodology}
\begin{enumerate}
    \item \textbf{20-day lag}: Accounts for the longer time from infection to death:
    \begin{itemize}
        \item Days 0--5: Incubation period
        \item Days 5--10: Symptom development
        \item Days 10--15: Hospitalization (for severe cases)
        \item Days 15--20: Death (for fatal cases)
    \end{itemize}
    \item \textbf{Daily data}: Used daily death counts from NYT data
    \item \textbf{Inclusion of zeros}: Days with zero deaths were included in the calculation
\end{enumerate}

\subsubsection{Output}
\begin{itemize}
    \item Counties analyzed: 3,218
    \item Total calculations: 1,899,366
    \item Output files: \texttt{CaseStudies/p\_ID\_analysis/}
\end{itemize}

\subsection{CFR: Case Fatality Rate}

\subsubsection{Formula}
\begin{equation}
    \text{CFR} = \frac{\text{Total Deaths (cumulative)}}{\text{Total Cases (cumulative)}}
\end{equation}

\subsubsection{Methodology}
\begin{enumerate}
    \item \textbf{End-of-period values}: Used cumulative cases and deaths as of December 31, 2021
    \item \textbf{Simple ratio}: No time lag applied (standard epidemiological definition)
\end{enumerate}

\subsubsection{Output}
\begin{itemize}
    \item Counties analyzed: 3,140
    \item Total U.S. cases: 52,568,947
    \item Total U.S. deaths: 781,929
    \item Output files: \texttt{CaseStudies/CFR\_analysis/}
\end{itemize}

\section{Validation: Carson City, NV (FIPS 32510)}

The automated methodology was validated against manual calculations for Carson City, Nevada:

\begin{table}[h]
\centering
\begin{tabular}{lccc}
\toprule
\textbf{Parameter} & \textbf{Manual Calculation} & \textbf{Automated} & \textbf{Match} \\
\midrule
$p_{IH}$ (mean) & 10.60\% & 10.60\% & $\checkmark$ \\
$p_{ID}$ (mean) & 0.1742\% & 0.1742\% & $\checkmark$ \\
CFR & 1.84\% & 1.84\% & $\checkmark$ \\
\bottomrule
\end{tabular}
\caption{Validation of automated calculations against manual Carson City analysis}
\end{table}

\section{File Structure}

\subsection{Scripts}
\begin{itemize}
    \item \texttt{extract\_hospitalization\_data.py}: Extracts county-level hospitalization CSVs
    \item \texttt{calculate\_pih\_all\_counties\_complete.py}: Calculates $p_{IH}$ for all counties
    \item \texttt{calculate\_pid\_all\_counties\_complete.py}: Calculates $p_{ID}$ for all counties
    \item \texttt{calculate\_cfr\_all\_counties\_complete.py}: Calculates CFR for all counties
\end{itemize}

\subsection{Output Directories}
\begin{itemize}
    \item \texttt{CaseStudies/hospitalization\_data/by\_state/}: County hospitalization CSVs (2,446 files)
    \item \texttt{CaseStudies/p\_IH\_analysis/}: $p_{IH}$ results and summaries
    \item \texttt{CaseStudies/p\_ID\_analysis/}: $p_{ID}$ results and summaries
    \item \texttt{CaseStudies/CFR\_analysis/}: CFR results and summaries
\end{itemize}

\section{Key Assumptions and Limitations}

\begin{enumerate}
    \item \textbf{Recovery period}: Assumed 14-day recovery for active case calculation
    \item \textbf{Hospitalization lag}: Assumed 14-day lag from infection to hospitalization
    \item \textbf{Death lag}: Assumed 20-day lag from infection to death
    \item \textbf{Data completeness}: Some counties may have incomplete reporting
    \item \textbf{Hospital aggregation}: Multiple hospitals per county were summed
    \item \textbf{Date range}: Analysis limited to March 2020 -- December 31, 2021
\end{enumerate}

\end{document}

